\cleardoublepage
\chapter{Metodyka badań}
\label{cha:MetodykaBadan}

W~niniejszym rozdziale omówiono metodykę realizacji badań. 
Przedstawiono scenariusze testowe, zastosowane konfiguracje eksperymentalne, 
parametry pomiarowe oraz narzędzia służące do rejestrowania czasu wykonania, obciążenia procesora i~wykorzystania pamięci RAM.

\section{Scenariusze testowe}
\label{sec:ScenariuszeTestowe}

Eksperymenty przeprowadzono według różnych konfiguracji parametrów, których celem jest zebranie informacji potrzebnych do przeprowadzenia analizy.

\subsection*{Zakres badanych konfiguracji}

W badaniach uwzględniono cztery algorytmy:
\begin{itemize}
    \item Quick Sort
    \item Merge Sort
    \item Bucket Sort
    \item Sample Sort
\end{itemize}

Dla każdego algorytmu wykonano pomiary jego wariantu sekwencyjnego oraz wariantu równoległego. 
Wariant sekwencyjny stanowił punkt odniesienia i był zawsze uruchamiany na jednym rdzeniu.

Testy przeprowadzono dla następujących rozmiarów danych zarówno dla liczb całkowitych, jak i zmiennoprzecinkowych:
\begin{itemize}
    \item 1~000 elemntów
    \item 10~000 elemntów
    \item 100~000 elemntów
    \item 1~000~000 elementów
\end{itemize}

Wykorzystano dwanaście wariantów zbiorów wejściowych, obejmujących:
\begin{itemize}
    \item dane losowe
    \item dane z duplikatami
    \item dane częściowo posortowane (20\%, 40\%, 60\% oraz 80\% elementów uporządkowanych)
\end{itemize}

\subsection*{Konfiguracja równoległości}

Testy równoległe wykonano dla dostępnych konfiguracji liczby jednostek wykonawczych udostępnianych przez środowisko sprzętowe. 
W przypadku algorytmów Quick Sort oraz Merge Sort liczba procesów była zależna od parametru głębokości równoległości:
\[
% \textit{max\_depth} = \log_2(\textit{cores}), 
\textit{max\_depth} = \left\lfloor \log_2(p) \right\rfloor, 
\]

gdzie p oznacza liczb˛e jednostek wykonawczych. W przypadku Bucket Sort oraz Sample Sort 
liczbę procesów ustawiano jako równą liczbie wykorzystanych jednostek wykonawczych.

\subsection*{Organizacja eksperymentów}

Każdy pomiar wykonywano wielokrotnie, a do analizy przyjmowano wartości uśrednione. 
Liczbę powtórzeń ustalono na 10 dla każdego wariantu eksperymentu, 
aby uzyskać możliwie dokładną wartość średnią z przeprowadzonych pomiarów. Przed właściwymi pomiarami wykonano dodatkowe uruchomienie,
którego wynik nie był uwzględniany w obliczeniach. Miało ono na celu przygotowanie środowiska do właściwych testów, jak również 
zprofilowanie uruchomionego algorytmu, dzięki czemu narzut pracy profilera nie wpływał na wyniki.

Przeprowadzone testy pozwoliły uzyskać zróżnicowany zestaw wyników, które umożliwiają analizę wpływu rozmiaru danych, 
liczby wykorzystywanych jednostek wykonawczych oraz charakterystyki zbioru wejściowego na czas sortowania, zużycie pamięci oraz efektywność zrównoleglenia.

\section{Wariant sekwencyjny jako punkt odniesienia}
\label{sec:WariantSekwencyjny}

Każdy z~badanych algorytmów, oprócz wersji równoległej, posiada również swój wariant sekwencyjny. 
Stanowił on punkt odniesienia dla przeprowadzanych po nim testów wersji równoległej, 
umożliwiając określenie, jak wariant równoległy wypada na tle odpowiadającego mu wariantu sekwencyjnego.

Uzyskane wyniki z wersji sekwencyjnej wykorzystywano następnie do wyznaczenia metryk efektywności zrównoleglenia:
\begin{itemize}
    \item przyspieszenia (\textit{speedup})
    \item efektywności (\textit{efficiency})
\end{itemize}

Wariant sekwencyjny jest wykonywany na pojedynczej jednostce wykonawczej, z wykorzystaniem identycznych danych wejściowych oraz tych samych 
metod pomiaru czasu sortowania, zużycia pamięci RAM i obciążenia procesora co w odpowiadającym mu wariancie równoległym. 
Jednakowe warunki umożliwiają skupienie analizy na rzeczywistych różnicach pomiędzy wariantem sekwencyjnym i~równoległym oraz uzyskiwanymi 
przez nie wynikami.

\section{Parametry pomiarów}
\label{sec:ParametryPomiarow}

Podczas każdego eksperymentu zbierano wyniki pozwalające na późniejszą analizę wydajności i~efektywności algorytmów 
oraz wykorzystania zasobów sprzętowych.

Podstawowym parametrem jest czas wykonania algorytmu, mierzony z wysoką rozdzielczością za pomocą funkcji \texttt{time.perf\_counter}. 
Dla serii powtórzeń zapisywano:
\begin{itemize}
    \item czas średni
    \item medianę
    \item odchylenie standardowe
\end{itemize}

Dodatkowo monitorowano zużycie zasobów systemowych:
\begin{itemize}
    \item średnie obciążenie procesora CPU w trakcie działania algorytmu
    \item średnie zużycie pamięci operacyjnej RAM
    \item maksymalne zużycie pamięci operacyjnej RAM
\end{itemize}

Na podstawie czasu sekwencyjnego $T_1$ oraz czasu równoległego $T_p$ 
wyznaczano metryki efektywności zrównoleglenia:
\begin{itemize}
    \item przyspieszenie (\textit{speedup}):
    \[
    S_p = \frac{T_1}{T_p},
    \]
    \item efektywność (\textit{efficiency}):
    \[
    E_p = \frac{S_p}{p} = \frac{T_1}{p \cdot T_p},
    \]
    gdzie $p$ oznacza liczbę jednostek wykonawczych.
\end{itemize}

Wartości \textit{speedup} oraz \textit{efficiency} dla wariantu sekwencyjnego są pomijane.

Wszystkie uzyskane wyniki zapisywano w bazie danych SQLite, 
co umożliwiło ich bezpieczne przechowywanie oraz wykorzystanie w dalszych etapach badań.

\section{Narzędzia do pomiaru czasu, CPU i pamięci}
\label{sec:NarzedziaPomiaru}

Do realizacji potrzebnych pomiarów wykorzystano zestaw narzędzi i bibliotek umożliwiających 
precyzyjne rejestrowanie czasu wykonania oraz monitorowanie wykorzystania zasobów systemowych w trakcie działania algorytmów.

\subsection*{Pomiar czasu}

Czas wykonania mierzono za pomocą funkcji \texttt{time.perf\_counter} z biblioteki standardowej Pythona~\cite{bib:PerfCounter}. 
Zapewnia ona dostęp do zegara o wysokiej rozdzielczości, przeznaczonego do pomiaru krótkich przedziałów czasu. 
Czas wykonania wyznaczano jako różnica wartości licznika zarejestrowanych przed rozpoczęciem i po zakończeniu
badanego fragmentu kodu. Dzięki wysokiej rozdzielczości funkcja dobrze nadaje się do przeprowadzania pomiarów wydajnościowych. 
W implementacji CPython wykorzystywany zegar jest monotoniczny, dzięki czemu jego wartość nie może cofać się w czasie.

\subsection*{Monitorowanie CPU i pamięci RAM}

Zużycie procesora oraz pamięci operacyjnej monitorowano z wykorzystaniem biblioteki \texttt{psutil}~\cite{bib:Psutil}. 
Pomiary realizowano w osobnym wątku, który w ustalonych odstępach czasu zapisanych w \texttt{sample\_interval} 
odczytywał informacje o wykorzystanych zasobach.
Dla zbiorów danych o rozmiarze nieprzekraczającym 10~000 odstęp wynosił 0,01 sekundy, natomiast dla pozostałych zbiorów 0,05 sekundy.
Podczas każdego pomiaru odczytywano:
\begin{itemize}
    \item łączne obciążenie CPU procesu głównego oraz procesów potomnych
    \item zużycie pamięci RAM na podstawie wartości RSS
\end{itemize}

Dzięki uwzględnieniu procesów potomnych możliwe było poprawne oszacowanie wykorzystywanych zasobów przez algorytmy równoległe.

\subsection*{Profilowanie wydajności}

Do analizy szczegółowej wydajności wykorzystano dwa narzędzia:

\begin{itemize}
    \item \texttt{cProfile}~\cite{bib:cProfile} -- wbudowany profiler Pythona, 
    umożliwiający pomiar liczby wywołań oraz czasu wykonania poszczególnych funkcji
    
    \item \texttt{py-spy}~\cite{bib:PySpy} -- zewnętrzny profiler próbkujący umożliwiający analizę 
    działania procesu głównego i podprocesów oraz generowanie wykresów typu \textit{flamegraph}
\end{itemize}

Profilowanie nie wpływało na wartości raportowane w głównych wynikach, ponieważ pierwsze uruchomienie, 
na którym działały profilery, było pomijane przy obliczaniu statystyk.

\subsection*{Organizacja pomiaru}

Każdy pojedynczy przebieg algorytmu uruchamiano w osobnym procesie roboczym, co zapewniało izolację pomiaru oraz 
umożliwiało zabezpieczenie testów w~przypadku awarii lub zawieszeniu badanego algorytmu.
Wprowadzono limit czasowy ustawiony na 15 minut, po przekroczeniu którego proces był przerywany, a dany test otrzymywał status \texttt{TIMEOUT}.

Zebrane próbki CPU i RAM agregowano do wartości średnich oraz maksymalnych, 
a następnie wraz z~wynikami pomiaru czasu oraz wartościami \textit{speedup} i~\textit{efficiency}